% ======================================================================
% 9. EFFECTS AND MITIGATION  (budget: 2 pages)
% Structure: damage -> controls by layer -> measured effect in our lab.
% Every control gets: what it stops, what it costs, what it does NOT stop.
% ======================================================================
\section{Effects and Mitigation}
\label{sec:mitigation}

\subsection{Damage model}
% TODO
% - PtH converts one compromised workstation into the privileges of every
%   privileged account that touched it. Quantify: one host, N cached admins.
% - The compounding effect over time: each hop adds hashes.
% - Historical anchors: NotPetya and similar campaigns as the reference case.

\subsection{Credential protection}
% TODO
% - Credential Guard / LSA protection (PPL): stops remote extraction of the
%   hash; explicitly state what it does not do (a local admin can still
%   abuse a live session).
% - Protected Users group and its consequences (no NTLM, no delegation, no
%   RC4) -- and the operational pain it causes.
% - Disable WDigest; restrict cached credentials.

\subsection{Protocol hardening}
% TODO
% - Disable or restrict NTLM (audit first, then block per service); this is
%   the only control that removes the PtH path outright.
% - Kerberos: remove RC4, require AES, enforce FAST/armoring -- this removes
%   the PtK path.
% - SMB signing and encryption: state precisely that they address relay and
%   tampering, not hash possession.
% - EPA / channel binding for the HTTP-facing cases.

\subsection{Architectural controls}
% TODO
% - Tiered administration, PAW, LAPS for local accounts, JIT elevation,
%   gMSA/service accounts. Frame as "reduce how many hashes are reachable
%   from a workstation", not as a checklist.

\subsection{Measured effectiveness in the lab}
% TODO -- the strongest section of the essay
% - Table: control applied -> scenario A/B/C outcome -> cost/limitation.
% - Include one control that did not help, and explain why. That single
%   negative result demonstrates understanding.
% Figure: fig-mitigation-matrix.
